% Options for packages loaded elsewhere
\PassOptionsToPackage{unicode}{hyperref}
\PassOptionsToPackage{hyphens}{url}
%
\documentclass[
]{article}
\usepackage{amsmath,amssymb}
\usepackage{iftex}
\ifPDFTeX
  \usepackage[T1]{fontenc}
  \usepackage[utf8]{inputenc}
  \usepackage{textcomp} % provide euro and other symbols
\else % if luatex or xetex
  \usepackage{unicode-math} % this also loads fontspec
  \defaultfontfeatures{Scale=MatchLowercase}
  \defaultfontfeatures[\rmfamily]{Ligatures=TeX,Scale=1}
\fi
\usepackage{lmodern}
\ifPDFTeX\else
  % xetex/luatex font selection
\fi
% Use upquote if available, for straight quotes in verbatim environments
\IfFileExists{upquote.sty}{\usepackage{upquote}}{}
\IfFileExists{microtype.sty}{% use microtype if available
  \usepackage[]{microtype}
  \UseMicrotypeSet[protrusion]{basicmath} % disable protrusion for tt fonts
}{}
\makeatletter
\@ifundefined{KOMAClassName}{% if non-KOMA class
  \IfFileExists{parskip.sty}{%
    \usepackage{parskip}
  }{% else
    \setlength{\parindent}{0pt}
    \setlength{\parskip}{6pt plus 2pt minus 1pt}}
}{% if KOMA class
  \KOMAoptions{parskip=half}}
\makeatother
\usepackage{xcolor}
\setlength{\emergencystretch}{3em} % prevent overfull lines
\providecommand{\tightlist}{%
  \setlength{\itemsep}{0pt}\setlength{\parskip}{0pt}}
\setcounter{secnumdepth}{-\maxdimen} % remove section numbering
% definitions for citeproc citations
\NewDocumentCommand\citeproctext{}{}
\NewDocumentCommand\citeproc{mm}{%
  \begingroup\def\citeproctext{#2}\cite{#1}\endgroup}
\makeatletter
 % allow citations to break across lines
 \let\@cite@ofmt\@firstofone
 % avoid brackets around text for \cite:
 \def\@biblabel#1{}
 \def\@cite#1#2{{#1\if@tempswa , #2\fi}}
\makeatother
\newlength{\cslhangindent}
\setlength{\cslhangindent}{1.5em}
\newlength{\csllabelwidth}
\setlength{\csllabelwidth}{3em}
\newenvironment{CSLReferences}[2] % #1 hanging-indent, #2 entry-spacing
 {\begin{list}{}{%
  \setlength{\itemindent}{0pt}
  \setlength{\leftmargin}{0pt}
  \setlength{\parsep}{0pt}
  % turn on hanging indent if param 1 is 1
  \ifodd #1
   \setlength{\leftmargin}{\cslhangindent}
   \setlength{\itemindent}{-1\cslhangindent}
  \fi
  % set entry spacing
  \setlength{\itemsep}{#2\baselineskip}}}
 {\end{list}}
\usepackage{calc}
\newcommand{\CSLBlock}[1]{\hfill\break\parbox[t]{\linewidth}{\strut\ignorespaces#1\strut}}
\newcommand{\CSLLeftMargin}[1]{\parbox[t]{\csllabelwidth}{\strut#1\strut}}
\newcommand{\CSLRightInline}[1]{\parbox[t]{\linewidth - \csllabelwidth}{\strut#1\strut}}
\newcommand{\CSLIndent}[1]{\hspace{\cslhangindent}#1}
\ifLuaTeX
  \usepackage{selnolig}  % disable illegal ligatures
\fi
\usepackage{bookmark}
\IfFileExists{xurl.sty}{\usepackage{xurl}}{} % add URL line breaks if available
\urlstyle{same}
\hypersetup{
  pdftitle={The Prompt Is Not the Author},
  pdfauthor={Researcher},
  pdfkeywords={prompting, scholarly
contribution, authorship, provenance, AI disclosure, prompt
programming, human-AI writing},
  hidelinks,
  pdfcreator={LaTeX via pandoc}}

\title{The Prompt Is Not the Author}
\usepackage{etoolbox}
\makeatletter
\providecommand{\subtitle}[1]{% add subtitle to \maketitle
  \apptocmd{\@title}{\par {\large #1 \par}}{}{}
}
\makeatother
\subtitle{Contribution, Control, and the Missing Scholarly Object in
AI-Assisted Research}
\author{Researcher}
\date{2026-08-18}

\begin{document}
\maketitle
\begin{abstract}
Institutions increasingly ask researchers to disclose the magnitude or
significance of generative-AI use, while technical research increasingly
treats model-facing instructions, constraints, templates, and chains as
programmable components of systems. These practices collide around a
unit that neither side handles well: the prompt. One response is to
claim that prompts are simply a new form of authorship. That claim is
too strong. Copyright doctrine shows why detailed instructions can
contain substantial intellectual labor without supplying sufficient
control over a model's expressive realization. At the same time,
programming systems such as LMQL and DSPy, prompt-chain interfaces,
provenance systems, and human-AI writing traces show why treating
prompts as mere conversational input is too weak. This paper argues for
a narrower position: consequential prompting should be evaluated as a
role-relative scholarly contribution when it functions as specification,
method, configuration, orchestration, selection, or an epistemic
instrument, without assuming that these roles collapse into expressive
authorship. The required unit of analysis is therefore not the final
AI-generated sentence and not a scalar amount of ``AI use,'' but a
contribution topology that records what operations were performed, where
causal leverage occurred, what degree of control was exercised, under
what execution envelope the intervention was calibrated, and how the
resulting claim can be traced. This account also sets limits on
provenance. A full chat transcript is neither necessary nor sufficient
for evaluating contribution, and audit requirements can become reactive,
changing the inquiry they purport to document. The paper concludes with
an admission test for prompt-based scholarly objects and a research
program for evaluating prompt-forward dissertations without reducing
them either to conventional authorship or to machine contamination.
\end{abstract}

\emph{Working Paper · AI-augmented research process · Evidence, prompts,
and making history preserved.}

\section{1. The wrong unit}\label{the-wrong-unit}

A researcher submits a chapter produced through repeated interaction
with a generative model. A reviewer asks, ``Did AI write this?'' The
question sounds precise because the chapter ends in sentences, and
sentences are familiar evidence of writing. Yet the question can
collapse radically different production histories. One researcher may
type a three-word request, accept five generated pages, and make minor
edits. Another may construct a source set, specify comparisons and
exclusions, run dozens of adversarial variants, reject all machine
prose, and retain only a conceptual distinction that emerged through the
interaction. A third may build a prompt program whose constraints and
control flow become part of a reproducible experiment. The final
artifact can make these workflows look similar even when their causal
and intellectual structures are not.

This problem is not solved by asking a slightly more quantitative
version of the same question: ``How much AI was involved?'' Georgia
Tech's current policy, for example, makes the magnitude and significance
of AI use, together with disciplinary norms, relevant to disclosure
(Georgia Institute of Technology 2026). Those are reasonable governance
predicates, but magnitude does not by itself tell us what kind of work
was done. Recent work on measuring human contribution in AI-assisted
generation makes the same point from another direction. Xie and
colleagues propose an information-theoretic measure of how much
information human input contributes to AI-assisted output rather than
inferring contribution solely from the final text (Xie et al. 2026).
Their intervention is useful precisely because raw interaction volume
and human contribution are not equivalent. A short input can exert
strong conditional influence; a long input can have little. But even
that improvement leaves open a second distinction: informational
leverage is not intellectual merit. A three-word instruction can be
causally decisive and scholarly trivial.

The central claim of this paper is therefore deliberately narrower than
``prompts are authorship.'' A prompt, prompt program, selection
procedure, or generative loop can become a scholarly contribution when
it performs a consequential role in inquiry that can be specified,
tested, bounded, and traced. That contribution may be methodological,
conceptual, computational, curatorial, or expressive. These categories
should not be collapsed. The U.S. Copyright Office's 2025 report on AI
copyrightability is especially useful here because it pressures the
strongest version of prompt authorship: detailed prompting may contain
substantial human intellectual work while still failing to provide
sufficient control over the expressive details of generated output (U.S.
Copyright Office 2025). Rather than treating this as a defeat,
scholarship should take it as a distinction it badly needs.

The relevant unit is not the origin of the last sentence. It is the
structure of contribution around the artifact.

\section{2. Authorship, contribution, and control are different
thresholds}\label{authorship-contribution-and-control-are-different-thresholds}

The word \emph{authorship} currently carries too many jobs. It can mean
legal ownership, expressive origination, intellectual contribution,
responsibility for claims, eligibility for a byline, or satisfaction of
a degree requirement. Generative systems expose the looseness because
they separate operations that conventional writing often bundled
together. One person can specify an argument without realizing every
sentence. Another can select and arrange generated material without
specifying its initial wording. A model can supply sentences without
accepting responsibility for them. A researcher can reject every
generated passage while allowing the interaction to redirect the
argument.

CRediT already offers institutional precedent for separating
contribution from authorship. Its taxonomy distinguishes
conceptualization, methodology, software, validation, investigation,
supervision, writing, and other roles while explicitly declining to
determine who qualifies as an author (NISO, n.d.). That choice matters.
It shows that scholarship can describe the work people performed before
deciding what higher-order status should follow. Prompt-mediated
research needs the same discipline. Before asking whether a prompt makes
someone ``the author,'' we should ask what operation the prompt
performed.

Copyright supplies an important contrary case. The U.S. Copyright Office
distinguishes human instructions from human authorship when the
generative system determines expressive realization. It also recognizes
that sufficiently creative human selection, coordination, arrangement,
or modification of generated material can itself contain protectable
authorship (U.S. Copyright Office 2025). Two consequences follow. First,
intellectual sophistication in a prompt does not guarantee expressive
control. Second, post-generation operations can sometimes be authorship
even when the initial generation was not. Prompt-forward scholarship
should preserve both findings instead of flattening them into a slogan.

This yields at least four thresholds that should be kept separate:

\begin{enumerate}
\def\labelenumi{\arabic{enumi}.}
\tightlist
\item
  \textbf{Causal contribution:} did an intervention materially change
  system behavior or the artifact?
\item
  \textbf{Scholarly contribution:} did the intervention perform a
  consequential intellectual or methodological role in the research?
\item
  \textbf{Expressive control:} did the researcher determine the relevant
  expressive elements strongly enough for the applicable authorship
  standard?
\item
  \textbf{Responsibility:} who accepts and can answer for the resulting
  scholarly claim?
\end{enumerate}

Ablation makes the distinction vivid. Generative Agents demonstrates
through ablations that architectural components such as observation,
planning, and reflection contribute to the believability of agent
behavior (Park et al. 2023). Removing a component and observing
degradation is evidence that the component is load-bearing. But it does
not by itself establish who deserves credit for the component, whether
the component is novel, or whether its creator authored the system's
outputs. Causal necessity is not scholarly merit. It is one evidentiary
layer.

This is also why the familiar director or conductor analogy has to be
used carefully. Orchestration can be consequential without being
identical to writing. CRediT already gives names such as methodology,
supervision, software, and project administration to contributions that
coordinate work without pretending that every form of coordination is
textual authorship (NISO, n.d.). Copyright, meanwhile, distinguishes
coordinating the production process from creatively arranging the
resulting expression. Prompt-forward scholarship needs both kinds of
orchestration in its vocabulary.

\section{3. Prompt is already too broad a technical
category}\label{prompt-is-already-too-broad-a-technical-category}

The scholarly debate often speaks of ``the prompt'' as though a chat
request, a templated instruction, a constrained query program, and a
recursive agent workflow were variations of one object. Technical
practice already undermines that assumption.

LMQL explicitly generalizes prompting from free text into language-model
programming by combining natural-language fragments with scripting,
constraints, and control flow (Beurer-Kellner, Fischer, and Vechev
2023). DSPy goes further in a different direction: it abstracts
language-model pipelines as parameterized modules and computational
graphs and compiles them against a metric, deliberately moving away from
hand-written prompt templates as the sole locus of control (Khattab et
al. 2023). DSPy Assertions turns constraints into executable checks that
can trigger refinement when a model output violates specified conditions
(Singhvi et al. 2023). AI Chains and PromptChainer similarly make
multi-step prompt workflows explicit, allowing outputs to become inputs
to later stages while exposing intermediate states for inspection and
debugging (Wu, Terry, and Cai 2022; Wu et al. 2022).

These systems make a simple but consequential point: natural-language
surface form does not determine computational role. A string can be
conversational input in one setting, versioned configuration in another,
and an element of a larger program in a third. The Generative Agents
repository makes this materially visible by storing model-facing prompts
as assets under the simulation engine rather than treating them as
disposable chat residue. The associated paper, by contrast, describes
the architecture primarily through higher-level concepts such as memory,
reflection, planning, and retrieval (Park et al. 2023). The prompt can
therefore be load-bearing implementation without being the paper's
preferred conceptual unit.

This suggests a rule for scholarly classification: \textbf{classify a
model-facing artifact by its use, execution context, and relation to the
research claim, not by the fact that it is made of natural-language
text.} The same sentence can be a protocol, a configuration value, a
research direction, an interview question, or a model instruction. Those
uses are not identical merely because the characters are identical.

The implication is not that every prompt should be elevated into a new
publication genre. On the contrary, prompt exceptionalism is another
category error. Most prompts are routine. Most lines of code are routine
too. The relevant question is when a model-facing artifact becomes
sufficiently load-bearing, reusable, discriminating, or epistemically
consequential that omitting it prevents the research from being
understood, tested, or reproduced.

\section{4. From scalar attribution to contribution
topology}\label{from-scalar-attribution-to-contribution-topology}

If ``how much AI?'' is too coarse, what should replace it? The answer
should not be a maximally detailed ontology of every click. A useful
representation must preserve distinctions that change an epistemic or
institutional judgment and discard distinctions that do not.

A starting point is to treat contribution as a topology with at least
five dimensions: \textbf{role, leverage, control, calibration, and
provenance}.

\subsection{4.1 Role}\label{role}

Role asks \emph{what kind of scholarly work the intervention performs}.
CRediT offers a ready vocabulary at a relatively high level:
conceptualization, methodology, software, validation, writing,
supervision, and so forth (NISO, n.d.). Prompt-forward work may require
lower-level operational verbs beneath those roles: specify, constrain,
retrieve, branch, compare, reject, select, verify, recurse, orchestrate,
regenerate. But those verbs should be admitted only when they preserve a
decision-relevant distinction. The goal is not a taxonomy of every
interaction event; it is a vocabulary that prevents ``AI use'' from
swallowing materially different research activities.

This immediately reframes disclosure. ``AI was used for brainstorming''
is weak because \emph{brainstorming} can hide workflows with different
roles and different consequences. CoAuthor demonstrates how fine-grained
writing interactions can be recorded across requests, suggestions,
acceptance, dismissal, and revision (Lee, Liang, and Yang 2022). Humanly
similarly starts from the premise that final text alone cannot reveal
how a document was produced and records writing and in-platform AI
assistance as process evidence (S. Zhu et al. 2026). These systems make
it possible, at least in principle, to derive disclosure from operations
rather than from broad self-selected labels.

\subsection{4.2 Leverage}\label{leverage}

Leverage asks \emph{how much changing this intervention changes the
result}. Xie et al.'s information-theoretic approach is one way to
formalize a related idea: human input can contribute varying amounts of
information to an AI-assisted output (Xie et al. 2026). Technical
systems supply another family of tests. Ablation asks whether the system
loses capability when a component is removed. Benchmarking asks how much
performance changes relative to a baseline. Constraint systems ask
whether violations decrease when the intervention is present.

But leverage must not be mistaken for merit. A hard-coded constant can
be indispensable. A short prohibition can dramatically change output. A
clever conceptual distinction can exert only modest immediate control.
The purpose of leverage is to establish causal position, not to award
scholarly credit automatically.

\subsection{4.3 Control}\label{control}

Control asks \emph{what aspects of the outcome the researcher actually
determines}. This is where copyright doctrine is most useful as a
limiting case. Detailed instructions can narrow desired content without
determining expressive realization (U.S. Copyright Office 2025).
Technical constraint systems sharpen the distinction further. A prose
instruction such as ``do not invent citations'' is a soft constraint:
the model may violate it. An LMQL constraint or a validator that rejects
prohibited outputs is a different mechanism because the desired
condition has been moved into executable machinery (Beurer-Kellner,
Fischer, and Vechev 2023). DSPy Assertions makes a similar distinction
operational by checking constraints and using failures in refinement
loops (Singhvi et al. 2023).

Control is therefore not adequately proxied by prompt length,
specificity, or apparent sophistication. A thousand-word prompt can
exert less reliable control than a short programmatic condition. If the
scholarly argument depends on control, control should be measured at the
level where it is exercised.

\subsection{4.4 Calibration}\label{calibration}

Calibration asks \emph{under what execution conditions the claimed
operation works}. POSIX demonstrates that model behavior can be
sensitive to intent-preserving variations in prompt wording (Chatterjee
et al. 2024). PromptBench similarly treats prompts as one component in a
broader evaluation environment involving tasks, models, prompt
construction, adversarial variations, and evaluation protocols (K. Zhu
et al. 2023). These results make the phrase ``battle-hardened prompt''
relational. A prompt tested repeatedly on one model may be robust only
to that model and context.

For prompt-based scholarly instruments, the operative unit should
therefore be an \textbf{execution envelope}: prompt or prompt program,
model and version, relevant parameters, supplied context or sources,
task, evaluation procedure, and stopping or acceptance conditions. Exact
textual replication may be impossible as proprietary models change. In
that case, prompt-forward scholarship should predeclare what form of
functional replication is expected: equivalent constraint satisfaction,
comparable behavior distributions, preservation of a methodological
distinction, or reproduction of a research conclusion.

\subsection{4.5 Provenance}\label{provenance}

Provenance asks \emph{how the relevant artifact and claim were
produced}. W3C PROV is useful because it refuses to equate provenance
with a raw stream. It represents entities, activities, agents, and
relations such as use, generation, derivation, association, and
influence (Moreau and Missier 2013). That is already closer to what
AI-assisted scholarship needs than a binary ``AI-generated'' label.

Yet provenance alone does not assign responsibility. A perfect
production graph can show that a passage was generated from a particular
prompt, selected after several trials, and edited into a chapter. It
still cannot answer who should defend the claim or whether the selection
constituted sufficient scholarly contribution. Provenance is descriptive
infrastructure; adjudication is a normative layer added on top.

The contribution topology therefore should not end in a single score. It
should expose enough structure for different institutions to ask their
own questions without pretending they are asking the same one.

\section{5. Provenance should not mean
transcript}\label{provenance-should-not-mean-transcript}

The most tempting response to uncertainty is to keep everything. If
final text cannot reveal the process, preserve the entire conversation
and let future evaluators inspect it. This appears transparent, but it
confuses event capture with useful provenance.

PROV's central insight is that provenance is modeled through selected
entities, activities, agents, and relations (Moreau and Missier 2013). A
complete raw log can be one evidentiary source for such a graph, but it
is not the graph itself. In AI-assisted scholarship, the consequential
structure may include source selection, model choice, prompt
construction, generation, rejection, verification, arrangement, and
acceptance of responsibility. A transcript contains some of those events
mixed with irrelevancies and may omit others entirely.

Human-AI writing systems make this design problem visible. HaLLMark
captures writer-LLM interaction and visualizes provenance so writers can
communicate AI use to readers and publishers (Hoque et al. 2023).
DraftMarks uses familiar material metaphors such as revision traces and
marks for AI-originated material to surface writing process in the final
artifact (Siddiqui et al. 2025). Humanly packages captured process into
a sealed certificate rather than asking readers to interpret a raw
session log (S. Zhu et al. 2026). These systems disagree in
representation, but they converge on a premise: \textbf{process evidence
must be designed for inspection}.

A prompt receipt should therefore be understood as an interface to
contribution-relevant provenance, not a screenshot dump. A useful
receipt might identify the execution envelope, the load-bearing
instructions, representative failed trials, acceptance or rejection
operations, verification steps, and the relation between those events
and the final claim. The right minimum is an empirical question and
probably differs by audience. A reader evaluating transparency, a
reviewer evaluating reproducibility, and an integrity officer evaluating
a disputed claim do not need identical traces.

This proposal must also be placed under pressure from ordinary
research-record norms. Guidance for laboratory notebooks has long
emphasized sufficiently complete records to support validation,
reporting, reproducibility, and credit. Failed experiments are not
automatically private residue. This weakens any claim that failed
prompts are uniquely entitled to disappear. The real difference may lie
elsewhere: AI conversations can be unusually voluminous,
conversationally intimate, and easy to repurpose as evidence about the
researcher's intent or dependence. The design problem is therefore not
simply ``retain less.'' It is how to combine completeness, access
control, purpose limitation, and useful abstraction.

\section{6. When provenance becomes
performative}\label{when-provenance-becomes-performative}

There is a second danger in treating process traces as straightforward
truth. Once actors know that a trace will be evaluated, the trace can
change the practice that produces it.

Espeland and Sauder call attention to \emph{reactivity}: public measures
and evaluations can reorganize the people and institutions they measure
(Espeland and Sauder 2007). Their subject is rankings, not AI prompting,
so the extension must remain a hypothesis. But the mechanism is directly
testable. If researchers know that prompt histories may be inspected for
evidence of independent authorship, they may begin asking safer
questions, avoiding strong machine-generated objections, or narrating
their own agency more conspicuously inside the log. The resulting record
could become more polished and less naturalistic.

This is not an argument against provenance. Laboratory notebooks are
also written under norms of possible inspection. It is an argument for
separating two properties that are too easily conflated: \textbf{event
completeness} and \textbf{behavioral naturalism}. More capture can
improve reconstruction while also increasing reactivity. The optimal
point should be measured rather than assumed.

The problem is especially acute for provenance visualization. DraftMarks
deliberately makes process legible through material metaphors of
revision (Siddiqui et al. 2025). Such interfaces may help readers reason
about production, but they can also smuggle a moral economy of visible
labor into authorship judgments. Revision marks, keystrokes, and prompt
volume can become proof-of-work signals even though effort and
intellectual contribution are not the same. A polished trace may be
interpreted as virtuous; a concise high-leverage intervention may look
suspiciously easy.

This opens a concrete research program. Researchers can manipulate audit
visibility while holding tasks constant and measure whether prompt
diversity, adversarial exploration, or willingness to entertain
counter-normative hypotheses changes. They can vary the visual rhetoric
of provenance while holding underlying contribution constant and measure
perceived ownership or authorship. These are empirical questions, not
philosophical decorations.

\section{7. Governance needs translation, not the wholesale import of
engineering
ontology}\label{governance-needs-translation-not-the-wholesale-import-of-engineering-ontology}

The technical literature already recognizes model-facing artifacts as
programs, constraints, templates, modules, or components. It would be
easy to conclude that governance is simply behind and should import the
ontology of AI engineering. That conclusion is also too strong.

Engineering and governance ask different questions. DSPy asks how to
represent and optimize language-model pipelines against task metrics
(Khattab et al. 2023). LMQL asks how to combine prompting with
scripting, constraints, and control flow (Beurer-Kellner, Fischer, and
Vechev 2023). Georgia Tech's policy asks when AI use must be disclosed,
what oversight is required, and how disciplinary norms should shape
responsible academic and research practice (Georgia Institute of
Technology 2026). The same artifact may be relevant to each regime
without receiving the same predicate. A highly effective prompt program
can be an excellent engineering component and also a form of AI
involvement that should be disclosed. Those judgments are not
contradictory.

The sharper problem is translation. Governance categories such as
\emph{magnitude}, \emph{significance}, and \emph{appropriate disclosure}
operate at a level that can hide distinctions technical practice makes
explicit. But technical ontology has blind spots of its own:
responsibility, authorship, consent, disciplinary norms, and the meaning
of a doctoral contribution are not optimization targets. What is needed
is a shared set of objects with different institutional predicates, not
the rule of one discipline over the other.

Georgia Tech's policy also forces a correction to one of the strongest
claims in the initial prompt-forward argument. The use of qualitative
predicates is not necessarily evidence that governance has failed to
define its object. Mature institutional standards often rely on
professional interpretation and precedent rather than numerical
thresholds. The relevant question is whether those standards are
adequately calibrated. Georgia Tech explicitly delegates part of the
judgment to disciplinary norms (Georgia Institute of Technology 2026).
For interdisciplinary prompt-forward work, this delegation can itself
become the problem: computing, HCI, literary studies, design, and
graduate education may classify the same operation differently.

This is why adjudicated exemplars matter. A general policy need not
contain a complete theory of prompting, but researchers need some way to
see how standards are applied. The evidentiary chain would include the
relevant AI interaction or operational artifact, the resulting scholarly
work, the disclosure, the institutional judgment, and the rationale.
Whether such exemplars exist across local units is an empirical
question. The absence of examples in a central policy is not proof of
arbitrary governance. It is a research edge.

\section{8. The prompt as epistemic
instrument}\label{the-prompt-as-epistemic-instrument}

The strongest positive claim in the field does not come from authorship
at all. It comes from the philosophy of scientific instruments.

Davis Baird argues that scientific knowledge is not exhausted by
propositions and that material instruments can themselves embody
epistemically significant know-how (Baird 2004). The analogy to
prompting must be handled carefully because Baird's force partly comes
from moving beyond language-centered epistemology, while a prompt is
made of language. Yet model-facing instructions can occupy a strange
hybrid position: they are linguistic artifacts whose important property
may be neither their propositional content nor the prose they produce,
but the repeatable operation they enable in a generative system.

This suggests the possibility of \textbf{operative scholarly writing}:
symbolic artifacts whose scholarly content is partly demonstrated by
what they reliably make a system do. The category would include only a
subset of prompts. An ordinary request for prose would not qualify. A
versioned prompt program that operationalizes a comparison, survives
specified perturbations, exposes its boundary conditions, and enables a
reproducible analytical procedure might.

The instrument analogy imposes obligations, not privileges. PromptBench
and POSIX show why calibration matters (K. Zhu et al. 2023; Chatterjee
et al. 2024). A prompt cannot be declared an instrument merely because
it produced a desired result once. The execution envelope must be
specified; sensitivity must be measured; failure modes must be known.
Baird also blocks the reverse error of attributing all
instrument-mediated knowledge to the user. An effective prompt system
inherits work from model builders, training data, frameworks, libraries,
prior prompt designers, and institutional infrastructure. The researcher
who operates the instrument does not automatically own the knowledge
embodied upstream.

A prompt-based scholarly object should therefore have to earn the
category through use.

\section{9. An admission test for prompt-based scholarly
objects}\label{an-admission-test-for-prompt-based-scholarly-objects}

The preceding distinctions support a practical admission test. A
model-facing artifact should be treated as a candidate scholarly object
when enough of the following are true:

\textbf{1. It specifies an operation.} The artifact does more than ask
for generic output. It encodes a comparison, constraint, transformation,
retrieval procedure, decision rule, or experimental structure that
matters to the research.

\textbf{2. It makes a demonstrable behavioral difference.} Removing,
replacing, or perturbing it changes a relevant capability or research
result relative to an appropriate baseline. This establishes leverage,
not merit.

\textbf{3. Its execution envelope is reported.} The model, version,
context, source set, relevant parameters, evaluation procedure, and
acceptance conditions are specified to the extent needed for functional
replication.

\textbf{4. Its boundaries are known.} The researcher can identify
failure cases, sensitivity, transfer limits, or contexts in which the
artifact should not be expected to work.

\textbf{5. Its construction has provenance.} The artifact's derivation
from experiments, prior prompts, source material, or collaborative work
can be traced without pretending that a raw transcript is itself the
scholarly object.

\textbf{6. It performs an identifiable scholarly role.} The artifact
contributes to conceptualization, methodology, software, validation,
writing, selection, orchestration, or another defensible role rather
than receiving credit merely for existing.

\textbf{7. The researcher accepts epistemic responsibility for the
claims made with it.} The system's participation does not substitute for
the researcher's obligation to verify sources, state limits, and answer
criticism.

This is intentionally not an authorship test. Some artifacts passing it
may also support expressive authorship; some may not. It is a test for
whether a prompt-based artifact belongs inside the scholarly record as
more than generic ``AI use.''

The test can itself be falsified. It may prove too elaborate, privilege
technical prompting, or fail to change committee judgments. The zettel
field points directly to experiments: compare free-form instructions
with structured programs; ablate prompts and prompt architectures
separately; evaluate workflow traces using scalar and graph
representations; test committee judgments before and after
contribution-role descriptions; measure prompt robustness across models
and versions; and vary audit visibility to test reactivity. The theory
should survive those experiments or change.

\section{10. What this means for a prompt-forward
dissertation}\label{what-this-means-for-a-prompt-forward-dissertation}

A prompt-forward dissertation should not try to win legitimacy by
claiming that every prompt is writing. That position is unnecessary and
vulnerable. It should instead show where the researcher's contribution
resides, preserve enough evidence to inspect it, and allow different
categories of contribution to remain different.

For a dissertation committee, the key question is not whether every word
was manually produced. Nor is it whether AI participation crossed an
abstract percentage threshold. The sharper question is which roles must
remain substantially the doctoral candidate's. CRediT cannot answer that
because it explicitly does not determine authorship or degree
requirements (NISO, n.d.). A graduate program must answer it
institutionally. But the contribution topology can make the answer
inspectable. A committee can ask whether the candidate originated the
research problem, constructed the method, designed or selected the
prompt programs, performed validation, diagnosed failures, made
consequential selections, integrated the argument, and can defend every
claim.

The answer may differ by chapter. A prompt can be method in one section,
configuration in a software artifact, documentary evidence in another,
and merely convenience tooling elsewhere. Placement should follow
function. If a paper's central claim depends on a prompt program,
exiling that artifact to an appendix can make the scholarship less
inspectable. If the prompt is routine and only documents ordinary
assistance, an appendix or concise disclosure may be enough.

This approach also makes room for a difficult possibility: a
prompt-forward dissertation can contain substantial AI-generated prose
while still requiring a large and demonstrable human research
contribution. Conversely, a dissertation can contain almost entirely
human-typed prose while outsourcing crucial conceptual judgment to a
model. Word-origin purity is not a reliable proxy for intellectual
independence.

The purpose of prompt receipts is therefore not to prove virtue. It is
to make consequential operations inspectable. The receipt should not say
``look how hard I worked.'' It should let another researcher or
committee member see what was specified, what changed, what failed, what
was selected, what was verified, and what remains the candidate's
responsibility.

\section{11. Limits and unresolved
territory}\label{limits-and-unresolved-territory}

Several limits prevent this account from becoming a new universal
ontology.

First, the evidence does not establish that current AI-disclosure policy
systematically penalizes high-impact prompting. Georgia Tech's policy
requires attention to magnitude, significance, and disciplinary norms,
but disclosure relevance is not condemnation (Georgia Institute of
Technology 2026). Claims that institutions treat prompt capability as
contamination require case evidence, committee interviews, or controlled
studies. The stronger pharmakon formulation should therefore remain a
hypothesis about institutional framing, not a description of policy
fact.

Second, the evidence does not establish that full prompt histories are
routinely demanded in misconduct investigations. That procedural claim
requires an institutional receipt. The paper instead establishes a more
modest point: if process histories are used for attribution, their
evidentiary design matters.

Third, contribution topology may become too complicated for policy. A
representation that faithfully records every node and edge of a
recursive human-AI workflow can defeat its own purpose. W3C PROV
supports abstraction, but the right compression level is claim-relative
and context-relative (Moreau and Missier 2013). Any practical standard
will have to trade descriptive fidelity against interpretability,
privacy, and reactivity.

Fourth, the account risks privileging what can be formalized. LMQL and
DSPy make technical prompt operations inspectable, but literary, design,
ethnographic, and tacit practices may resist the same object language
(Beurer-Kellner, Fischer, and Vechev 2023; Khattab et al. 2023). A
plural operational vocabulary may be necessary.

Fifth, the instrument analogy remains incomplete. Prompts are symbolic,
model-relative, and often unstable under software updates. Baird's
material ``thing knowledge'' cannot simply be transferred to them (Baird
2004). The category of operative scholarly writing is therefore a
proposal generated by the collision, not an established lineage.

These limits are not cleanup. They locate the next work.

\section{12. Conclusion}\label{conclusion}

The prompt becomes intellectually interesting at the point where
familiar categories stop agreeing about it. Copyright can treat a
detailed instruction as insufficient expressive control. Programming
research can treat model-facing instructions and constraints as
components of executable systems. Human-AI writing research can show
that final text underdetermines the process that produced it. Provenance
standards can represent the process without determining responsibility.
Institutional policies can require disclosure of significant AI
involvement without supplying a theory of prompt contribution.

The mistake is to force these regimes into one verdict. A prompt need
not be ``the author'' to be scholarly work. It need not be copyrightable
expression to function as method. It need not contain final prose to
redirect an argument. It need not be long to exert high causal leverage.
And high leverage alone does not make it intellectually important.

What scholarship needs is a way to keep those facts visible at once. The
proposed contribution topology begins with five questions: \textbf{What
role did the intervention perform? How much leverage did it exert? What
did the researcher actually control? Under what execution envelope was
the operation calibrated? What provenance is necessary to reconstruct
and adjudicate the contribution?} Those questions do not eliminate
authorship. They stop asking one overloaded noun to carry the entire
factory.

The deepest research question therefore survives, but in altered form.
Not: \emph{Should prompts count as writing?} Not even: \emph{Should
prompts be disclosed?} The question that remains is more operational and
more difficult:

\textbf{What kind of scholarly object is an instruction, constraint, or
generative procedure that changes what a machine can reliably do, and
what evidence is sufficient to show the researcher's contribution to
that change?}

That is where a prompt-forward dissertation can make a defensible wager.

\section{Appendix A --- Assembly
Instrument}\label{appendix-a-assembly-instrument}

The exact assembly instrument used to package this working paper and its
zettel field is preserved in
\texttt{\_PROMPTS/SLIPCASE\_ASSEMBLY\_PROMPT.txt}. The prompt is treated
as evidence of the publication procedure, not as evidence for the
paper's substantive claims.

\section{Appendix B --- Making History}\label{appendix-b-making-history}

The package distinguishes material \textbf{provided} by the researcher,
material \textbf{preserved} from the active research context, sources
\textbf{retrieved or checked} during packaging, artifacts
\textbf{derived} by compilation, and claims that remain
\textbf{unverified}. The complete record appears in
\texttt{000\_\_MAKING\_HISTORY.txt} and the paper-specific
making-history file.

\section{Appendix C --- Replication
Path}\label{appendix-c-replication-path}

Open \texttt{index.html} for the sendable offline research desk. The
immutable zettel payloads remain available as root \texttt{.txt} cards
and exact \texttt{\_MD/} mirrors. \texttt{ZETTELS.jsonl},
\texttt{\_SLIPCASE/NODES.jsonl}, and \texttt{\_SLIPCASE/RELATIONS.jsonl}
preserve machine-readable state. \texttt{000\_\_RETURN\_PATH.txt}
records the checkpoint identity and rejoin phrase;
\texttt{000\_\_REBUILD.txt} records the rebuild path.

\section*{References}\label{bibliography}
\addcontentsline{toc}{section}{References}

\phantomsection\label{refs}
\begin{CSLReferences}{1}{0}
\bibitem[\citeproctext]{ref-baird2004thing}
Baird, Davis. 2004. \emph{Thing Knowledge: A Philosophy of Scientific
Instruments}. University of California Press.

\bibitem[\citeproctext]{ref-beurerkellner2023lmql}
Beurer-Kellner, Luca, Marc Fischer, and Martin Vechev. 2023.
{``Prompting Is Programming: A Query Language for Large Language
Models.''} In \emph{Proceedings of the ACM SIGPLAN Conference on
Programming Language Design and Implementation}, 1946--69.
\url{https://arxiv.org/abs/2212.06094}.

\bibitem[\citeproctext]{ref-chatterjee2024posix}
Chatterjee, Anwoy, H. S. V. N. S. Kowndinya Renduchintala, Sumit Bhatia,
and Tanmoy Chakraborty. 2024. {``POSIX: A Prompt Sensitivity Index for
Large Language Models.''} In \emph{Findings of the Association for
Computational Linguistics: EMNLP 2024}, 14550--65.
\url{https://aclanthology.org/2024.findings-emnlp.852/}.

\bibitem[\citeproctext]{ref-espeland2007rankings}
Espeland, Wendy Nelson, and Michael Sauder. 2007. {``Rankings and
Reactivity: How Public Measures Recreate Social Worlds.''}
\emph{American Journal of Sociology} 113 (1): 1--40.
\url{https://doi.org/10.1086/517897}.

\bibitem[\citeproctext]{ref-gatech2026ai}
Georgia Institute of Technology. 2026. {``Artificial Intelligence (AI)
in Academic \& Research Contexts.''}
\url{https://policylibrary.gatech.edu/artificial-intelligence-ai-academic-research-contexts}.

\bibitem[\citeproctext]{ref-hoque2023hallmark}
Hoque, Md Naimul, Tasfia Mashiat, Bhavya Ghai, Cecilia Shelton, Fanny
Chevalier, Kari Kraus, and Niklas Elmqvist. 2023. {``The HaLLMark
Effect: Supporting Provenance and Transparent Use of Large Language
Models in Writing with Interactive Visualization.''} \emph{arXiv
Preprint arXiv:2311.13057}. \url{https://arxiv.org/abs/2311.13057}.

\bibitem[\citeproctext]{ref-khattab2023dspy}
Khattab, Omar, Arnav Singhvi, Paridhi Maheshwari, Zhiyuan Zhang, Keshav
Santhanam, Sri Vardhamanan, Saiful Haq, et al. 2023. {``DSPy: Compiling
Declarative Language Model Calls into Self-Improving Pipelines.''}
\emph{arXiv Preprint arXiv:2310.03714}.
\url{https://arxiv.org/abs/2310.03714}.

\bibitem[\citeproctext]{ref-lee2022coauthor}
Lee, Mina, Percy Liang, and Qian Yang. 2022. {``CoAuthor: Designing a
Human-AI Collaborative Writing Dataset for Exploring Language Model
Capabilities.''} In \emph{Proceedings of the 2022 CHI Conference on
Human Factors in Computing Systems}.
\url{https://doi.org/10.1145/3491102.3502030}.

\bibitem[\citeproctext]{ref-moreau2013prov}
Moreau, Luc, and Paolo Missier, eds. 2013. {``PROV-DM: The PROV Data
Model.''} World Wide Web Consortium.
\url{https://www.w3.org/TR/prov-dm/}.

\bibitem[\citeproctext]{ref-niso_credit}
NISO. n.d. {``CRediT Role Descriptors.''}
\url{https://credit.niso.org/contributor-roles-defined/}.

\bibitem[\citeproctext]{ref-park2023generative}
Park, Joon Sung, Joseph C. O'Brien, Carrie J. Cai, Meredith Ringel
Morris, Percy Liang, and Michael S. Bernstein. 2023. {``Generative
Agents: Interactive Simulacra of Human Behavior.''} In \emph{Proceedings
of the 36th Annual ACM Symposium on User Interface Software and
Technology}. \url{https://arxiv.org/abs/2304.03442}.

\bibitem[\citeproctext]{ref-siddiqui2025draftmarks}
Siddiqui, Momin N., Nikki Nasseri, Adam Coscia, Roy Pea, and Hari
Subramonyam. 2025. {``DraftMarks: Enhancing Transparency in Human-AI
Co-Writing Through Interactive Skeuomorphic Process Traces.''}
\emph{arXiv Preprint arXiv:2509.23505}.
\url{https://arxiv.org/abs/2509.23505}.

\bibitem[\citeproctext]{ref-singhvi2023assertions}
Singhvi, Arnav, Manish Shetty, Shangyin Tan, Christopher Potts, Koushik
Sen, Matei Zaharia, and Omar Khattab. 2023. {``DSPy Assertions:
Computational Constraints for Self-Refining Language Model Pipelines.''}
\emph{arXiv Preprint arXiv:2312.13382}.
\url{https://arxiv.org/abs/2312.13382}.

\bibitem[\citeproctext]{ref-usco2025copyrightai}
U.S. Copyright Office. 2025. {``Copyright and Artificial Intelligence,
Part 2: Copyrightability.''} U.S. Copyright Office.
\url{https://www.copyright.gov/ai/Copyright-and-Artificial-Intelligence-Part-2-Copyrightability-Report.pdf}.

\bibitem[\citeproctext]{ref-wu2022promptchainer}
Wu, Tongshuang, Ellen Jiang, Aaron Donsbach, Jeff Gray, Alejandra
Molina, Michael Terry, and Carrie J. Cai. 2022. {``PromptChainer:
Chaining Large Language Model Prompts Through Visual Programming.''} In
\emph{CHI Conference on Human Factors in Computing Systems Extended
Abstracts}. \url{https://doi.org/10.1145/3491101.3519729}.

\bibitem[\citeproctext]{ref-wu2022aichains}
Wu, Tongshuang, Michael Terry, and Carrie J. Cai. 2022. {``AI Chains:
Transparent and Controllable Human-AI Interaction by Chaining Large
Language Model Prompts.''} In \emph{Proceedings of the 2022 CHI
Conference on Human Factors in Computing Systems}.
\url{https://doi.org/10.1145/3491102.3517582}.

\bibitem[\citeproctext]{ref-xie2026humancontribution}
Xie, Yueqi, Tao Qi, Jingwei Yi, Xiyuan Yang, Ryan Whalen, Junming Huang,
Qian Ding, Yu Xie, Xing Xie, and Fangzhao Wu. 2026. {``Measuring Human
Contribution in AI-Assisted Content Generation.''} In \emph{Proceedings
of the 64th Annual Meeting of the Association for Computational
Linguistics}, 6168--90.

\bibitem[\citeproctext]{ref-zhu2023promptbench}
Zhu, Kaijie, Qinlin Zhao, Hao Chen, Jindong Wang, and Xing Xie. 2023.
{``PromptBench: A Unified Library for Evaluation of Large Language
Models.''} \emph{arXiv Preprint arXiv:2312.07910}.
\url{https://arxiv.org/abs/2312.07910}.

\bibitem[\citeproctext]{ref-zhu2026humanly}
Zhu, Shenzhe, Haoqian Zhang, Xu Yang, Jingyu Tang, Yi Nian, Xiaoxue Du,
Shu Yang, Alex Pentland, Joachim Baumann, and Jiaxin Pei. 2026.
{``Humanly: A Configurable and Traceable Environment for Human-AI
Collaborative Writing.''} \emph{arXiv Preprint arXiv:2607.21758}.
\url{https://arxiv.org/abs/2607.21758}.

\end{CSLReferences}

\end{document}
